\section{Remarks and open questions}\label{sec:discussion}

\begin{remark}[What the results rest on]
For $p\ge2$, for $p\le1$, and for the minorants of \cref{thm:low}, the proofs use the
definition of rank-one convexity, one-variable calculus, and the fact that polyaffine functions
are null Lagrangians. The quasiconvexity of $F$ for $1<p<2$ rests in addition on the
quasiconvexity of the Burkholder function $B_p$, which is the content of Iwaniec's conjecture
in the form of~\cite[Conjecture~1.1]{GuerraKristensen2022} and is stated
in~\cite[Corollary~2]{SharpHI}. Some input of this kind is needed: the positive part of $B_p$
satisfies~\Hyp\ and is nonnegative, so an unconditional form of \cref{cor:low}\,(a) would
contain its quasiconvexity.
\end{remark}

\begin{remark}[The role of the hypotheses]
Isotropy and homogeneity enter at every step, through the profile and the function $h$ of one
variable. Nonnegativity is used in three places for $p>2$: in Case III, where the positive part
of the minorant is used; in Case I, when $\psi(1)\le0$; and at the ends. It is necessary there
(\cref{thm:iff}\,(a)). For $1<p<2$ it is used in \cref{prop:low}, and we do not know whether
\cref{cor:low}\,(a) holds for integrands that change sign.
\end{remark}

\begin{remark}[Open questions]
\begin{enumerate}[label=\textup{(\roman*)}]
\item For $p>2$ a function with~\Hyp\ that takes negative values is not polyconvex. Whether it
  is quasiconvex is open in general; see \cref{sec:intro-related}.
\item The arguments do not bear on Morrey's problem for integrands on $\M$ that are not
  isotropic or not homogeneous.
\end{enumerate}
\end{remark}
